\hypertarget{future-work}{%
\chapter{Future Work}\label{future-work}}

This chapter develops the future-work programme that arises directly from the limitations documented in Section 6.6 of the discussion and from the operational considerations discussed in Section 6.4. Five lines of investigation are described, each motivated by either a specific finding or limitation of the evaluation or by the operational requirements of deployment in a regulatory food-safety authority.

\hypertarget{enhancements-to-ai-models}{%
\section{Enhancements to AI Models}\label{enhancements-to-ai-models}}

The most immediate technical priority is the architectural mitigation of the label-leak failure mode documented in Section 5.5 and Section 6.3. The proposed mitigation is the introduction of an output-validation stage between the structured generation step and the audit-log persistence step: any emitted base-term code is resolved against the FoodEx2 Appendix B taxonomy at generation time, the label associated with the code is replaced by the official taxonomy label, and any code that does not appear in the retrieved candidate set triggers an automatic regeneration request with an explicit instruction to choose from the retrieval context. This validation can be implemented as a constraint within the existing generation function (Section 4.5.3) without architectural restructuring and is expected to eliminate label leak in the 11.8\% of cases where it currently occurs under the GPT-4o configuration and in the 23.5\% of cases where it occurs under the Claude configuration (Section 5.7.2).

The second technical priority is the redesign of the human-review trigger to address the 4.2\% recall problem documented in Section 5.6. The proposed redesign replaces the single-signal confidence-based trigger with a multi-signal trigger that incorporates four complementary inputs. The first is the rerank score of the top-1 candidate, which flags classifications where the top candidate scores below a configurable threshold. The second is the rerank-score spread among the top-5 candidates, which flags classifications where the top-5 candidates are too closely clustered to support a confident choice. The third is the agreement between the VLM's preliminary code and the final classifier's output, which flags cases where the RAG layer overrode the VLM on the rationale that disagreement signals lower confidence. The fourth is the presence of a documented coverage gap, which flags classifications that fall on the small list of base terms used as fallbacks when no specific term is available. The empirical calibration of these thresholds against the 34-image evaluation, and against a larger validation set as it becomes available, is itself a research subtask.

The third technical priority, building on the open-source extension already executed in the present work, is the broadening of the cross-VLM evaluation across additional open-source model families and parameter scales. The proprietary cross-VLM comparison reported in Section 5.7.2 established convergence between GPT-4o and Claude at 0.0\% bare-VLM accuracy and 29.4\% end-to-end system accuracy. The first open-source extension, against LLaVA-NeXT-Mistral-7B, was completed for the present submission and demonstrated that the system remains operable but that end-to-end accuracy collapses to 2.9\% at this model scale, attributable to a generic-fallback hallucination regime in which the open-source VLM emits the same A0FQJ "Mixed dish with meat and vegetables" code in 33 of 34 cases, contaminating the RAG query uniformly across the test set. The natural continuation is the evaluation of additional contemporary 7B-parameter open-source VLMs from other model families, which would test whether the LLaVA failure profile is idiosyncratic to that family or reproduces across the contemporary 7B open-source VLM class. The broader continuation comprises three sub-priorities: a parameter-scale sweep evaluating open-source VLMs at 13B, 34B, and 72B parameter counts where available, to localise the parameter-count threshold at which open-source models approach proprietary accuracy; an instruction-tuning study in which an open-source model is fine-tuned on FoodEx2 documentation and labelled food imagery, to test whether task-specific adaptation of the VLM layer can close the gap at the smaller parameter counts that fit local deployment hardware (this fine-tuning is positioned as a complement to, not a replacement for, the RAG-grounded generation layer that the present evaluation establishes as a necessary component of the system regardless of VLM choice, the architectural intent being a hybrid configuration in which a task-tuned open-source VLM provides a stronger preliminary signal for the same downstream retrieval and reranking pipeline, in line with recent evidence that hybrid RAG-and-fine-tuning pipelines outperform either component in isolation on domain-specific structured tasks); and the construction of a cost-and-accuracy trade-off curve across the full proprietary--open-source landscape, to inform deployment decisions in regulatory contexts where the proprietary-API dependency is an architectural concern. The substitutability of the VLM layer was a design-time property confirmed by the three working implementations; the broader cross-VLM evaluations would close the empirical case across the full open-source VLM landscape.

A fourth methodological priority, complementary to the architectural changes above, is the enrichment of the evaluation framework with metrics tailored to the hierarchical structure of the FoodEx2 taxonomy. The exact-match and family-level accuracy metrics used in the present work are simple, auditable, and reproducible from the per-image artefacts, but they do not credit partial correctness within the deeper levels of the FoodEx2 hierarchy beyond the first three characters of the base term. A hierarchical F1 (hF1) score computed against the full taxonomy path would provide a finer-grained measure of system performance, distinguishing for example between an error that drops the classification one taxonomy level from the ground truth and one that crosses to an unrelated branch. Adoption of hF1 alongside the existing accuracy measures, with the existing measures retained as the headline figures for cross-paper comparability, is identified as the recommended refinement for any subsequent evaluation campaign, particularly the larger-scale offline evaluation discussed in Section 8.5.

\hypertarget{expansion-of-foodex2-coverage}{%
\section{Expansion of FoodEx2 Coverage}\label{expansion-of-foodex2-coverage}}

The cultural-coverage limitations of the FoodEx2 taxonomy documented in Section 6.3 motivate a programme of contribution to the EFSA maintenance cycle. The seven dish categories surfaced as coverage gaps in the present evaluation, acarajé, brigadeiro, francesinha, pastel, rabanada, cuscuz nordestino, and a generic Brazilian-style mixed-meat-and-rice dish, are candidates for formal proposal to EFSA as new base terms in subsequent maintenance releases. The format of the EFSA maintenance reports (EFSA Supporting Publications, 2019--2025) is well established and the proposal mechanism is openly documented. The IDRISK2 system's audit logs would provide quantitative evidence of the operational need for these terms by recording each instance in which the existing taxonomy required a fallback to a generic parent code.

Beyond the specific proposals, the evaluation suggests a more systematic methodology in which IDRISK2 audit data is mined for clusters of classifications that converge on the same fallback base term across visually distinct dishes, a pattern that signals an underlying coverage gap of operational rather than incidental significance. This would constitute a feedback loop in which deployed instances of the system inform the maintenance cycle of the underlying taxonomy, aligning the regulatory ontology more closely with the operational reality of inspectors in non-European jurisdictions.

\hypertarget{integration-with-external-systems}{%
\section{Integration with External Systems}\label{integration-with-external-systems}}

The present implementation operates as a standalone classification service and does not integrate with the regulatory information systems used by ASAE, SIT, FVST, or other competent food-safety authorities. Operational deployment in any of these contexts will require integration with the laboratory information management systems (LIMS) in which sample metadata is maintained, with the inspection-workflow systems in which classifications are recorded as part of the inspection record, and ultimately with the EFSA data-collection infrastructure in which national authorities report their classifications upstream. Each of these integrations is a separate engineering project, but each is supported by the present system's emission of structured FoodEx2 outputs alongside audit records in standard JSON, which can be mapped to the target systems' schemas without modification to the core classification pipeline.

A specific integration of significant value is the alignment of the IDRISK2 audit trail with the EFSA Data Collection Framework. If audit records can be exported in a format compatible with EFSA's submission tooling, the burden of manually entering FoodEx2 classifications into national reporting systems is reduced and the consistency of the data collected by EFSA from multiple national authorities is improved. The technical implementation of this export mapping is straightforward; the institutional and contractual arrangements under which it would operate are the substantive work.

\hypertarget{advanced-security-features}{%
\section{Advanced Security Features}\label{advanced-security-features}}

The closed deployment architecture documented in Section 4.6 provides foundational capabilities that support compliance with GDPR Article 22's data-residency, data-minimisation, and audit-traceability requirements, and that align with the technical documentation, human oversight, and post-market monitoring obligations of the EU AI Act for high-risk AI systems. Full certification of a high-risk AI deployment, however, would require an extensive conformity assessment, a quality management system aligned with the relevant ISO standards (notably ISO/IEC 42001), and external audits that go substantially beyond the architectural capabilities implemented in the present prototype. The enhancements proposed in this section should therefore be read as strengthening foundational compliance capabilities rather than completing the certification process; three further enhancements would strengthen the system's security posture for deployments handling higher volumes or more sensitive data.

First, the audit-log HMAC chain currently writes one entry per classification to a flat directory of JSON files. For deployments at higher volumes, this storage pattern becomes operationally awkward and the chain verification becomes computationally expensive. A migration to an append-only log structure with periodic Merkle-tree commitments to a tamper-evident external anchor would preserve the cryptographic verification property at scale while supporting efficient incremental verification. Second, the role-based access control framework currently authenticates users through a User dataclass instantiated by the calling code, without an external identity provider. Integration with an organisational SSO provider, with cryptographic verification of the user's identity at request time, would close the trust gap between the calling code and the access control mechanism. Third, deployments processing data subject to additional regulatory frameworks beyond GDPR (for example, data covered by national food-safety incident reporting obligations) would benefit from the integration of a privacy-preserving aggregation mechanism, such as the differentially private aggregation primitives discussed in Section 3.4.2, that would allow the system's classification statistics to be shared across deployments without exposing individual classification records.

\hypertarget{longitudinal-studies-and-real-world-deployment}{%
\section{Longitudinal Studies and Real-world Deployment}\label{longitudinal-studies-and-real-world-deployment}}

The empirical evaluation presented in Chapter 5 was conducted on a small dataset by a single researcher and does not exercise the system under the operational conditions of an actual inspection workflow. Three categories of longitudinal study are needed to establish the system's operational viability.

First, a larger-scale offline evaluation with multi-annotator ground truth would address the statistical-uncertainty, inter-annotator-agreement, and geographical-bias limitations documented in Section 6.6. A dataset of several hundred to several thousand images, geographically balanced across the full EFSA regulatory remit so that the accuracy figures reflect operational European conditions rather than the Portuguese-and-Brazilian-biased lower-bound estimate of the present evaluation, with each image annotated by at least two independent FoodEx2-trained annotators, would tighten the confidence intervals on the present accuracy metrics, establish the practical ceiling on classification accuracy imposed by inter-annotator disagreement, and surface failure modes that the 34-image dataset is too small to expose.

Second, a controlled field deployment with cooperating inspectors from one of the target regulatory authorities, for example, ASAE in Portugal, would expose the system to the operational conditions of an actual inspection workflow. The deployment would measure the system's effect on inspector workload, classification consistency, and case-throughput under realistic time pressure, and would surface usability defects and integration friction that an offline evaluation cannot detect. A participatory design methodology in which inspectors contribute to the refinement of the system's prompts, the calibration of the human-review trigger, and the format of the audit record would be required to produce a deployment that is operationally compatible with existing inspection practice rather than a research prototype imposed on inspectors from outside. The field deployment is moreover the principal evidence base for the human-review trigger redesign proposed in Section 8.1: by observing which classifications inspectors choose to override under operational time pressure and how often they accept the system's high-confidence outputs without independent verification, the multi-signal trigger thresholds (rerank score, score spread, VLM-classifier disagreement, coverage-gap flags) can be calibrated against actual operator decision patterns rather than against the offline-derived heuristics of the present implementation. Grounding trigger calibration in observed inspector behaviour is necessary because the EU AI Act Article 14 human-oversight obligation requires not only that a trigger exists but that it escalates the cases that human reviewers actually find decision-critical, a property the present 34-image offline evaluation cannot establish.

Third, a longitudinal study tracking the evolution of system performance as the underlying FoodEx2 taxonomy is updated through the EFSA maintenance cycle would establish the system's resilience to taxonomy change. The IDRISK2 architecture is designed to support taxonomy updates without retraining (the knowledge base is rebuilt from the new Appendix B XLSX and the system continues to operate on the updated index), but this design property has been validated only at deployment time, not across a sequence of real updates. A study running over multiple maintenance cycles would measure whether the system's accuracy is stable, improving, or degrading as the taxonomy evolves, and would identify the maintenance practices required to keep the deployment current with the regulatory state of the art.
